\section{练习}

\begin{question}\begin{drill}[运动员对传送带的功率]\label{em3:drill:1}
运动员腰部由轻绳经左侧的定滑轮连接悬挂质量$m_2$的重物，重物和运动员重心均保持静止。传送带在脚下以速率$v$向左运动，运动员对带的水平作用力向左。忽略绳轮损耗，求运动员对带做功的功率。
\end{drill}
\fig{fig:em3:drill:1:s}{绳拉力向左，运动员在带面上施加的水平力也向左}{\begin{tikzpicture}

\draw[world](0,0)rectangle(4.2,.3);\draw[flow](3,-.5)--(1.5,-.5)node[midway,below]{$v$};
\fill(2.4,2)circle(.13);\draw[world](2.4,1.87)--(2.4,.9) (2.4,.9)--(1.9,.3) (2.4,.9)--(2.9,.3);
\draw[world](-.3,1.2)circle(.2);\draw[world](2.4,1.4)--(-.3,1.4)arc(90:180:.2)--(-.5,-.7);
\draw[world](-.85,-.7)rectangle(-.15,-1.3);\node[left]at(-.85,-1){$m_2$};
\draw[guide](-.3,1.2)--(-.3,1.8);\draw[world](-.6,1.8)--(0,1.8);
\node[right]at(3,1.5){人体重心静止};
\end{tikzpicture}}
\end{question}
\answerlink{em3:drill:1}
\begin{question}\begin{drill}[偏置钉子的完整圆周条件]\label{em3:drill:2}
长$L$轻绳一端固定于$A$，$A$在水平轴$Ox$上方$L/2$。小球从$A$右侧与$A$等高的位置静止释放。钉子$P$在$x>0$处，绳下摆时被其阻挡。绳最大张力为$7mg$，求球能绕$P$完整一周而不断绳的$x$范围。
\end{drill}
\fig{fig:em3:drill:2:s}{题设：偏置钉子的完整圆周条件}{\begin{tikzpicture}[scale=.9]\coordinate(A)at(0,1.5);\coordinate(P)at(1.4,0);\draw[axis](0,0)--(4.5,0)node[right]{$x$};\draw[guide](A)--(0,0);\draw[world](A)--(3,1.5);\fill(3,1.5)circle(.08);\fill(A)circle(.06)node[left]{$A$};\fill(P)circle(.06)node[below]{$P$};\node[left]at(0,.7){$L/2$};\draw[guide](0,1.5)--(1.4,0)node[midway,left]{$AP$};\draw[dim](0,-.5)--(1.4,-.5)node[midway,below]{$x$};\end{tikzpicture}}
\end{question}
\answerlink{em3:drill:2}
\begin{question}\begin{drill}[绳松弛后击中钉子]\label{em3:drill:3}
长$\ell$的单摆悬于$O$。钉子$C$在$O$正下方距离$x<\ell$处。球从左侧不高于$O$的位置静止释放，经过最低点后绳绕钉，随后松弛，球作抛体运动并恰好击中$C$。求$x$的最小值。
\end{drill}
\fig{fig:em3:drill:3:s}{原摆长$\ell$不变，绕钉后的运动半径为$\ell-x$}{\begin{tikzpicture}

\coordinate(O)at(0,3);\coordinate(A)at(-2.95442,2.47906);\coordinate(C)at(0,1.32853);
\draw[world](O)--(A)node[midway,above]{$\ell$};\fill(O)circle(.05)node[above]{$O$};\fill(A)circle(.08)node[left]{释放};
\draw[guide](-3.3,3)--(.4,3);\draw[guide](O)--(0,0);\fill(C)circle(.05)node[right]{$C$};
\draw[dim](.7,3)--(.7,1.32853)node[midway,right]{$x$};\draw[dim](.7,1.32853)--(.7,0)node[midway,right]{$r=\ell-x$};
\draw[guide](0,0)arc(-90:65:1.32853);\fill(0,0)circle(.06)node[below]{最低点};
\draw[flow](-2.7,1.8)--(-2.1,1)node[left]{下摆};
\end{tikzpicture}}
\end{question}
\answerlink{em3:drill:3}
